% A.73: full root and DE proof/source/exact-certificate audits PASS.
\subsection{Every semidihedral geometric closure}

\begin{statement}{A.73}{Proved}{Both cusp assertions for every semidihedral two-group closure}
Let \(Y\to E\setminus\{O\}\) be an actual connected finite
\'{e}tale cover whose geometric Galois closure has group
\[
 G=SD_{4n}=\langle r,s\mid r^{2n}=s^2=1,
                              \ srs^{-1}=r^{n-1}\rangle,
 \qquad n\ge4\text{ a power of two}.                 \tag{A.73.1}
\]
Then both assertions of Task A.5 hold for \(Y\) at every
coefficient prime, every induced cusp packet, and every
permitted open derived subgroup. Thus the full fixed geometric
pro-\(\ell\) subgroup is its selected cusp inertia and every
nontrivial arithmetically normalized procyclic pro-\(\ell\)
subgroup is cuspidal.

The proper Galois closure has genus \(2n-1\), four cusps,
potential good Jacobian reduction of special \(2\)-rank zero
at two, and cusp differences killed by \(2n\). The proof
includes every actual intermediate. The universal question
A.45 remains open.
\end{statement}

\begin{proof}
Let \(C\to E\) be the smooth proper compactification of
the actual geometric Galois closure. Enlarge the number field
finitely to define it, its deck transformations, quotient maps
and marked points. At a place over two enlarge further so that
the original elliptic curve has good reduction and the stable
models and endomorphisms used below are defined. In particular
every local automorphism fixes the coefficient valuation ring.

\paragraph{Exact group structure and actual cusp inertia.}
Put \(k=n-1\). Every group element is uniquely \(r^a s^b\),
with \(a\bmod2n\) and \(b\in\{0,1\}\), and
\[
                (a,b)(c,d)=(a+k^b c,\ b+d\bmod2).
                                                        \tag{A.73.2}
\]
Here \(k^2\equiv1\pmod{2n}\). A pair with one rotation
and one outside element generates \(G\) precisely when
the rotation exponent is odd. A pair \(r^a s,r^b s\)
generates precisely when \(a-b\) is odd. Indeed their
quotient is \(r^{a-b}\); in the even-difference case
all available rotations have even exponent, since
\((r^a s)^2=r^{na}\). Two rotations cannot generate.

Direct multiplication gives
\[
 [r^a,r^b s]=r^{(2-n)a},\qquad
 [r^a s,r^b s]=r^{(2-n)(a-b)}.                    \tag{A.73.3}
\]
For a generating pair these exponents are twice odd integers,
because \(4\mid n\). The peripheral commutator in the
geometric fundamental group of the punctured elliptic curve
therefore gives the actual inertia subgroup
\[
             I=\langle r^2\rangle,
             \qquad |I|=n,\qquad G/I\simeq C_2^2.
                                                        \tag{A.73.4}
\]
It is normal, so the actual Galois cover has four cusps,
each of ramification index \(n\), and
\[
                       g(C)=2n-1.                 \tag{A.73.5}
\]
This follows from \(2g(C)-2=4(n-1)\). Also \(I=\Phi(G)\):
the quotient by \(I\) is elementary abelian, squares and
commutators lie in \(I\), and the square of \(r\) generates
\(I\). Explicitly, \(HI=G\) forces a subgroup \(H\) to
contain an odd rotation and an outside element, hence \(H=G\).

\paragraph{The actual old factor and the special rank dichotomy.}
The nilpotent ramified-cover theorem of Lehr--Matignon applies
to the actual \(C\to E\): the base has a good smooth model,
the sole branch section \(O\) is \'{e}tale over its valuation
ring, and ramification in generic characteristic zero is tame.
Its Theorem 6.13 gives potential good reduction of \(J_C\)
and a tree as the proper stable dual graph.
\cite{fghLehrMatignon2005,fghSaidi2007}
Only this conclusion is used; no marked-fibre slope theorem
is assumed.

The central involution is \(\sigma=r^n\). The actual
quotient \(C_0=C/\langle\sigma\rangle\) has dihedral
group \(D_{2n}\), since \(n-1\equiv-1\pmod n\).
Its four cusp indices are \(n/2\), and \(g(C_0)=n-1\).
The actual classification in A.70 and the reduction statements
of A.44/A.69 show that its Jacobian has potential good
reduction of special \(2\)-rank zero.

Let \(A\) be the anti-invariant isogeny factor for
\(\sigma\). Actual pullback and norm give
\[
                J_C\sim J_{C_0}\times A,
                \qquad \dim A=n.
                                                        \tag{A.73.6}
\]
On this nonzero factor \(r^n=-1\). Since
\(x^n+1=\Phi_{2n}(x)\), the actual deck endomorphism
therefore gives a unital field injection
\(\mathbf Q(\zeta_{2n})\hookrightarrow\operatorname{End}^0(A)\).
After good reduction this acts on the rational \'{e}tale
Tate space of the special \(2\)-divisible group. The local
field \(\mathbf Q_2(\zeta_{2n})\) has degree \(n\),
whereas that Tate space has dimension at most \(\dim A=n\).
Consequently \(A\) has special \(2\)-rank zero or is ordinary.
We exclude the ordinary alternative.

\paragraph{What ordinarity says about the actual stable components.}
Take the stable marked model of \((C,S)\), with the four
cusps marked, over the resulting complete valuation ring.
Its graph is a tree, and its good special Jacobian is the
product of the Jacobians of its positive-genus components.
Suppose \(A\) is ordinary. The invariant factor for
\(\sigma\) has rank zero. A pair of positive-genus components
exchanged by \(\sigma\) is impossible: its invariant and
anti-invariant factors are both isogenous to the same
component Jacobian. Components on which \(\sigma\) is
the identity are retained in the invariant factor.

For a component \(B\) fixed by \(\sigma\) with a
nontrivial induced involution, let \(B'=B/\sigma\).
Actual pullback--norm identities give \(f(B')=0\), and
its anti-invariant Prym is ordinary if nonzero. Write
\(r_B\) for the number of fixed points and \(d_P\) for
the different exponents. The smooth-curve formulas are
\[
 \begin{aligned}
 f(B)&=r_B-1,\\
 g(B)-g(B')&=g(B')-1+\frac12\sum_{P\ {\rm fixed}}d_P,
                   \qquad d_P\ge2.
 \end{aligned}                                         \tag{A.73.7}
\]
They are Deuring--Shafarevich and Riemann--Hurwitz on the
smooth component curves.
\cite[Corollary 1.8]{fghCrew1984}
\cite[Proposition 1.4]{fghYang2024}
There is at least one fixed point. A zero-dimensional Prym
would force \(g(B')=g(B)=0\). Otherwise ordinarity gives
\(f(B)=g(B)-g(B')\), and hence
\[
                    r_B=g(B')+\frac12\sum d_P.
\]
It follows that \(g(B')=0\) and every \(d_P=2\).
Thus every positive-genus anti-invariant component is
ordinary, with \(\sigma\) its genus-zero-quotient involution;
its entire Jacobian is anti-invariant. The genera of all such
components sum to \(n\).

An ordinary curve with this involution admits no order-four
automorphism whose square is \(\sigma\). Here is the precise
input, also proved in A.69. For an ordinary curve with a
faithful \(2\)-group action, subtracting twice
Deuring--Shafarevich from Riemann--Hurwitz makes each
nonnegative term \(d_P-2(e_P-1)\) vanish. The classical
different sum and \(G_0=G_1\) then give \(G_2=1\).
The coefficient of \(t^2\) injects \(G_1/G_2\) into the
additive residue field, so an order-four subgroup cannot
fix a point. A fourth root of \(\sigma\) induces an
involution of \(\mathbf P^1=B/\sigma\), which has one
fixed point in characteristic two. Its fibre must be
ramified: otherwise the square of its permutation on a
two-point fibre would be trivial while \(\sigma\) exchanges
those two points. Thus the fourth root fixes the unique
point above it, a contradiction.
\cite[\S2.2, equation (3) and Exercise 2.12(c)]{fghBouwWewers}

\paragraph{The only possible component orbit.}
Let \(B\) be a positive-genus anti-invariant component, and
let \(G_B\) be its stabilizer. It contains \(\sigma\).
Every rotation outside \(\{1,\sigma\}\) has a power
of order four with square \(\sigma\), so such a rotation
cannot belong to \(G_B\). Every outside element \(r^a s\)
with \(a\) odd has square \(\sigma\), so it too is
excluded. These exclusions do not presume faithful action
of the whole stabilizer: an element whose square acts as
the nontrivial involution necessarily acts with order four.
Therefore
\[
             G_B\cap\langle r\rangle=\langle\sigma\rangle,
             \qquad |G_B|\le4.                       \tag{A.73.8}
\]
The second assertion also follows from the index-two rotation
subgroup. The orbit of \(B\) has at least \(n\) components,
all with positive anti-invariant genus. Their total genus is
only \(n\). Thus there is exactly one full \(G\)-orbit of
\(n\) such components, all ordinary elliptic curves, each
with stabilizer of order four. There are no other
positive-genus anti-invariant components.

Put \(\mu=r^2\). It belongs to the normal inertia subgroup
\(I\), so it fixes all four generic cusp points individually.
For each ordinary elliptic component \(B\), (A.73.8) gives
\[
            |\langle\mu\rangle B|=n/2\ge2.
                                                        \tag{A.73.9}
\]
In particular \(\mu\) moves every component in this orbit.
On each such ordinary elliptic curve, (A.73.7) says that
\(\sigma\) has exactly two fixed points.

\paragraph{A smooth fixed point would lift to an individually fixed cusp.}
Neither of those two points can be a smooth point of the
whole special fibre. Suppose it were, and complete the
smooth local chart there as \(R[[t]]\). The local
automorphism \(\sigma\) fixes \(R\). Its restriction
to \(B\) is nontrivial, so \(\sigma(t)-t\) has nonzero
reduction in the residue power-series ring. This reduction
vanishes to a positive finite order at the fixed point.
Weierstrass preparation therefore writes the difference
as a unit times a positive-degree distinguished polynomial,
with no uniformizer factor. It gives a nonempty finite flat
horizontal fixed locus in this formal disc.
\cite[Corollary 2.16]{fghVonk2021}
After a finite extension one of its generic points is an
actual point of the smooth proper generic curve, specializing
to the chosen point; this is the usual smooth formal-disc
identification and properness.

Any generic point fixed by \(\sigma\) is a cusp: the
actual cover \(C\to E\) is \'{e}tale off those points,
so its deck action is free there. All such cusps are fixed
by \(\mu\). Their specializations must consequently be
\(\mu\)-fixed. A smooth point on \(B\) cannot be fixed
by an automorphism moving its unique component, contrary
to (A.73.9). This proves that both \(\sigma\)-fixed
points on each ordinary elliptic component are nodes.
It is an explicit finite-flat lifting argument, not an
assumption that every wild special fixed point lifts.

\paragraph{The root must be fixed by the whole group.}
Let \(\Gamma\) be the finite underlying dual tree of
the stable marked model. Its centre is invariant under
every automorphism: repeatedly delete all terminal vertices
until a single vertex or a single edge remains. Choose
this vertex as root, or the midpoint of this edge. In the
second case an abstract barycentric subdivision makes it
a vertex; this does not alter the geometric components,
their genera, their nodes or their stability. Thus the
root is fixed by \emph{all} of \(G\), not just by
\(\mu\).

All \(n\) ordinary elliptic vertices belong to one
\(G\)-orbit, so they have the same positive distance
from this root. In particular none lies below another.
The root and each such vertex are fixed by \(\sigma\),
so their unique joining path is fixed pointwise. The node
on that elliptic component toward the root is one of its
two \(\sigma\)-fixed nodes. The other is a child node.
Its descendant subtree is nonempty and contains no other
ordinary elliptic vertex, because all those vertices have
the same depth.

This child subtree contains no marked cusp. Indeed \(\mu\)
fixes every cusp vertex and the chosen \(G\)-fixed root.
An automorphism of a finite tree fixing the endpoints of a
path fixes that path pointwise. A root-to-cusp path through
the moving elliptic vertex would contradict (A.73.9).
This is why a root fixed only by \(\mu\) would be
insufficient: its two elliptic orbits would not necessarily
have the same depth.

An unmarked finite child subtree with just one attachment
to its parent cannot consist entirely of rational components.
If it had \(v\) such vertices, its total valence within
the full graph would be \(2v-1\), whereas marked stability
requires at least \(3v\). Thus it contains a positive-genus
component. The child subtrees from the \(n\) equidistant
elliptic vertices are pairwise disjoint. None contains another
one of those elliptic vertices. Summing genera therefore gives
\[
                  g(C)\ge n+n=2n,
\]
contrary to (A.73.5). This excludes the ordinary anti-invariant
factor and proves that the actual good special Jacobian has
\(2\)-rank zero.

\paragraph{Actual cusp divisors and every intermediate.}
The quotient \(E'=C/I\) is an actual connected unramified
degree-four cover of \(E\), since the whole cusp inertia
has been killed. With an origin chosen above \(O\), it is
an elliptic isogeny with kernel \(C_2^2=E'[2]\).
For its four kernel points \(Q_i\), the actual map
\(h:C\to E'\) has one point \(P_i\) above each and
\(h^*Q_i=nP_i\). Pulling back a rational function with
divisor \(2(Q_i-Q_j)\) gives
\[
                    2n(P_i-P_j)\text{ principal}.
                                                        \tag{A.73.10}
\]
This uses the actual total ramification of the intermediate
map, not a substitute Jacobian factor.

For every proper subgroup \(H<G\), the cusp number of
\(D=C/H\) is \([G:HI]\ge2\), by the normal Frattini
property. Its genus is positive because it maps finitely
and nonconstantly to \(E\). The actual Jacobian norm
makes \(J_D\) an isogeny quotient of \(J_C\), hence it
has potential good reduction and special \(2\)-rank zero.
Norms of the functions in (A.73.10) kill every cusp
difference of \(D\) by \(2n\): pushforward of a point
over an algebraically closed field has coefficient one,
independently of ramification.

The torsion-divisor criterion in A.41 gives a rational
joining path to a distinct companion for every cusp of
\(D\). A.42 gives its full fixed inertia at every
coefficient prime, and A.59 gives every normalizer assertion
for its two-group closure. The quotient \(H=G\) is the
already settled original punctured elliptic curve.
Finite extensions are used for the algebraic data and
reduction only. Base-point paths are compared up to inner
conjugacy, and A.41--A.42's restriction argument handles
the original packets and permitted open derived groups.
No quotient-to-fixed-subgroup surjectivity is asserted.

The exact group and dimension checks are retained in
\path{code/a_semidihedral_two_power.py}. The tree and local
fixed-locus arguments above establish the unbounded
conclusion independently of that finite verification range.
\end{proof}
